\section{From a progression to five squareclasses}\label{sec:vertices}

By an affine change of coordinates the ten terms become \(\{\pm1,\pm3,\pm5,\pm7,\pm9\}\). Their root polynomial is even. Newton’s identities allow only a change of the constant coefficient, so the competing polynomial is also even. Its rational roots come in opposite pairs, with any zero occurring with even multiplicity. We may thus write them as \(\pm y_1,\ldots,\pm y_5\). The odd moments agree automatically, leaving
\begin{equation}\label{eq:four-even-moments}
 \sum_{i=1}^{5}y_i^{2j}=\sum_{a\in\{1,3,5,7,9\}}a^{2j}
 \qquad(1\leq j\leq4).
\end{equation}
No distinctness or nonvanishing assumption is imposed. Put
\[
 \mathcal A=\{1,9,25,49,81\},\qquad S_i=y_i^2,\qquad
 U=\prod_i S_i,\qquad U_0=945^2,
\]
and define
\begin{equation}\label{eq:H-polynomial}
 H(X)=X^4-165X^3+8778X^2-172810X+1057221.
\end{equation}
Newton's identities, applied through degree four, give
\begin{equation}\label{eq:split-quintic}
 \prod_{i=1}^{5}(X-S_i)=XH(X)-U,\qquad
 XH(X)-U_0=\prod_{a\in\mathcal A}(X-a).
\end{equation}
Thus the unfixed coefficient is itself a square, since
$U=(\prod_i y_i)^2$. We will show that the vertex belonging to one
specified squareclass must equal $25$; the second identity will then
recover all five vertices.

\begin{lemma}[Denominators]\label{lem:progression-denominators}
Write $y_i=n_i/D$, where $D>0$ and
$\gcd(n_1,\ldots,n_5,D)=1$. Every prime divisor of $D$ is congruent to
$1$ modulo $10$ and divides none of the $n_i$. All five numerators are
odd. In particular, the five fractions have the same reduced denominator,
and $S_i\equiv1\pmod8$ in $\Z_2$. If one $y_i$ is integral, then
$\{S_i\}_{i=1}^5=\mathcal A$ as multisets.
\end{lemma}
\begin{proof}
Set $M_j=\sum_{a\in\{1,3,5,7,9\}}a^{2j}$. Clearing denominators gives
\begin{equation}\label{eq:cleared-moments}
 \sum_i n_i^{2j}=M_jD^{2j}\qquad(1\leq j\leq4).
\end{equation}
If $2\mid D$, fourth powers modulo $16$ show that the number of odd
$n_i$ is divisible by $16$. That number lies between zero and five,
so every $n_i$ is even, contrary to primitivity. If $3\mid D$, use
sixth powers modulo $9$: a unit has sixth power one and a nonunit
has sixth power zero. The number of units is divisible by nine and
at most five, again contradicting primitivity.

The prime $5$ requires one further modulus. Suppose $5\mid D$.
Fourth powers modulo $5$ imply that the number of unit numerators
is zero or five. The first possibility contradicts primitivity,
so all five are units. Write $n_i^4=1+5a_i\pmod{25}$.
The fourth-power equation in \eqref{eq:cleared-moments} gives
$5+5\sum_i a_i\equiv0\pmod{25}$, or $\sum_i a_i\equiv-1\pmod5$.
Squaring these congruences now gives
\[
 \sum_i n_i^8\equiv5+10\sum_i a_i\equiv20\pmod{25},
\]
whereas the eighth-power equation requires zero. Thus $5\nmid D$.

Let $p>5$ divide $D$ and put $x_i=n_i^2\bmod p$. The four power
sums vanish. Since $1,2,3,4$ are invertible modulo $p$, Newton's
identities make the first four elementary symmetric functions vanish.
The root polynomial is therefore $X^5-e_5$. If $e_5=0$, every $n_i$
vanishes modulo $p$, which is impossible. Otherwise this polynomial
has five distinct nonzero roots in $\F_p$. Their ratios exhaust the
fifth roots of unity, so $5\mid p-1$. As $p$ is odd, $p\equiv1\pmod{10}$.
The roots are nonzero, proving also that $p\nmid n_i$ for every $i$.

We have proved that $D$ is odd. Now $M_2=9669\equiv5\pmod{16}$ and
$D^4\equiv1\pmod{16}$. Fourth powers count the odd numerators, forcing
all five to be odd. Their squares, divided by the odd square $D^2$,
are one modulo eight.

If some $y_i$ is integral, coprimality with $D$ gives $D=1$.
The equation $\sum_i n_i^2=165$ bounds every $|n_i|$ by
$\sqrt{165}<13$. If none belongs to $\{1,3,5,7,9\}$, all five odd
absolute values must equal $11$, contradicting the same equation.
A shared square is a common root of the two quintics in
\eqref{eq:split-quintic}; evaluating their difference gives $U=U_0$.
Their root multisets then agree.
\end{proof}

The restriction on denominator primes is sharp for these four equations
considered locally. If $p\equiv1\pmod{10}$ and $e\ge1$, put $D=p^e$.
Choose five nonzero residues $n_i$ whose squares are the fifth roots of
unity. Such choices exist because squaring is an automorphism of the
subgroup of order five. The four cleared equations hold modulo $p$.
Fix one numerator and lift the other four. The derivative matrix has
entries $2j n_i(n_i^2)^{j-1}$, for $1\le j\le4$. Its determinant is
a product of units times the Vandermonde determinant of four distinct
squares. Hensel's lemma therefore supplies a solution over $\Z_p$,
and all five resulting $y_i$ have valuation $-e$. Thus the global
argument cannot be replaced by a prohibition on negative valuations
at every prime. The dyadic condition, in contrast, places every rational
candidate in the single residue class used in Section~\ref{sec:marked}.

The denominator restriction gives affine reduction at all primes needed
below. It is important here that reduction retains multiplicity: a test
that only looked for five candidate roots would not exclude a repeated
vertex.

\begin{proposition}[Five prescribed classes]\label{prop:five-classes}
For a root $r$ of $H$, the vertices can be labeled by
$a\in\mathcal A$ so that
\begin{equation}\label{eq:five-squareclasses}
 \frac{S_a-r}{a-r}\in\Q(r)^{\times2}.
\end{equation}
Each label occurs exactly once; in particular, all vertices are nonzero
and pairwise distinct.
\end{proposition}
\begin{proof}
The certificate \cert{check_type18_product_classes.m}
gives exact assertions for the finite computations below. They use two-cover descent
for genus-two curves as implemented in Magma; see \cite{BruinStoll,magma}.
Every class-group bound used in this certificate is a Minkowski bound.

Let $d(U)=\operatorname{disc}_X(XH(X)-U)$. Explicitly,
\begin{align*}
 d(U)={}&3125U^4+108276547500U^3-1117199475459133506U^2\\
       &+533350028141359954915500U
         +2470517618419162907684321203125.
\end{align*}
The split quintic has square discriminant, the square of the
Vandermonde product. Consequently there are rational points with
parameters $U$ and $S_i$ on, respectively,
\[
 C_d:\quad V^2=Ud(U),\qquad C_H:\quad V^2=SH(S).
\]
For the second assertion use $S_iH(S_i)=U$.

First enumerate, for
\[
 p\in\{13,17,19,29,37,73,89,139,193\},
\]
all squares $u\in\F_p$ for which $XH(X)-u$ splits into square roots
with total multiplicity five. The only result is $u=U_0\bmod p$;
its roots are the five distinct elements of $\mathcal A\bmod p$.
None of these primes is $1\pmod{10}$, so Lemma~\ref{lem:progression-denominators}
applies. In particular even one such reduction rules out a zero or
repeated global vertex and shows $U\ne0$.

Here is the precise local comparison used to prune the descent classes.
For $f=d$ or $H$, the algebra of the two-cover map for $V^2=Xf(X)$
contains a rational factor at $X=0$. Represent a candidate by a polynomial
$w(X)$ and normalize it to $\widetilde w(X)=w(X)/w(0)$. The class of a
point with parameter $z$ has normalized representative $(z-X)/z$.
At a good test prime, factor $f$ into monic irreducibles $f_\nu$.
A necessary condition for compatibility with that parameter is
\begin{equation}\label{eq:resultant-filter}
 \operatorname{Res}\!\left(f_\nu,
       \widetilde w(X)\frac{z-X}{z}\right)\in\F_p^{\times2}
 \quad\text{for every }\nu.
\end{equation}
Indeed these resultants are norms of the relevant components, and
norms of squares are squares. If a representative has a pole or a
vanishing component modulo $p$, the certificate skips that comparison;
it does not discard the class. Thus this is a necessary local test,
with no assumption that norms detect all squares.

For $C_d$, descent initially gives five classes. Testing the sole
allowed parameter $U_0$ at $13$ and $19$ leaves two and then one,
namely the class represented by $U_0-X$. For $C_H$ there are initially
38 classes. Allowing any of the five seed residues, the successive
tests at $13,17,29,37,73,89,139,193$ leave
\[
 30,\ 22,\ 14,\ 13,\ 10,\ 8,\ 6,\ 5
\]
classes. The final set consists precisely of the classes $a-X$ for
$a\in\mathcal A$. At $89$, each seed residue is compatible with just
its own class. The five global vertices reduce to the five distinct
seed residues, so each global class occurs exactly once. Finally the
rational components $S_a$ and $a$ are squares. Their ratio removes
the possible common scalar in the cover comparison, giving
\eqref{eq:five-squareclasses} in the quartic component.
\end{proof}

The two auxiliary curves separate two roles of the same split quintic.
The discriminant curve records the compatibility of its five roots as
a collection; the vertex curve records an individual square root
parameter. The proof above uses the latter classes individually, with
the discriminant descent providing an additional check on the product
parameter. In particular, the passage to a labeled vertex does not
choose the smallest real root, impose an ordering on the $y_i$, or
assume that their signs agree with the seed. Labels are determined by
the local cover classes. This distinction lets the proof apply directly
to rational multisets arising from an even root polynomial, before
any chain decomposition has been made.

This assignment is stronger than membership in a union of five covers.
It forces a point on each labeled elliptic curve. The label $25$ will
suffice, and its dyadic residue is already fixed by the denominator lemma.

\section{The marked elliptic curve}\label{sec:marked}

Let $S=S_{25}$. We specify all fields and models used in the rank and
local arguments. Set
\begin{align}
 K&=\Q(h),&h^4-h^3-991h^2+13345h-19455&=0,\label{eq:quartic-K}\\
 r&=\frac{-h^3-23h^2+585h+5002}{73},&k&=25-r.\label{eq:r-and-k}
\end{align}
The element $r$ is a root of $H$, and
\[
 q(T)=\frac{H(T)}{T-r}=T^3+c_2T^2+c_1T+c_0,
\]\[
 c_2=r-165,\qquad c_1=r^2-165r+8778,\qquad
 c_0=r^3-165r^2+8778r-172810.
\]
Since $(S-r)/k$ and $H(S)=U/S$ are squares, $kq(S)$ is a square in
$K$. Therefore $S$ supplies a point
\begin{equation}\label{eq:marked-curve}
 P(S)=(kS,Y)\in E(K),\qquad
 E:\ Y^2=X^3+kc_2X^2+k^2c_1X+k^3c_0.
\end{equation}
The marked point is
\begin{equation}\label{eq:marked-point}
 P_0=(25k,189k),
\end{equation}
because $H(25)=189^2$. The following rank computation is unconditional;
its class-group input is established before the descent is used.

\subsection{A certified generating bound for the class group}

\begin{lemma}\label{lem:type18-class-group}
The cubic $q$ is irreducible over $K$. Its root field $M$ has degree
12, signature $(12,0)$, discriminant
\[
 D_M=5^6 11^9 1231^5 2417^5,
\]
and $\Cl(M)\simeq\Z/2\Z$. Prime ideals of norm at most $1000$ generate
$\Cl(M)$.
\end{lemma}
\begin{proof}
An absolute defining polynomial for $M$ is
\begin{align*}
 B(T)={}&T^{12}-3T^{11}-5625T^{10}+44481T^9+9843099T^8
 -99843252T^7\\
 &-5998855201T^6+67870752036T^5+854290494639T^4\\
 &-9750057898095T^3-2947764849192T^2
 +213431600705238T-420745260940011.
\end{align*}
The certificate \cert{type18_descent/rank.m} checks the signature and
maximal-order discriminant, irreducibility of $q$, and an isomorphism
$M\simeq K[T]/(q)$. This is also the field of an ordered pair of
distinct roots of the quartic \eqref{eq:quartic-K}.

We establish first that primes of norm at most $b=85\,000\,000$
generate. Put $X_*=2^{31}$, $\ell=\log X_*$, and
\[
 F(x)=\frac{\max(1-|x|/\ell,0)}{\cosh(x/2)},\qquad
 \Phi(s)=\int_{-\infty}^{\infty}F(x)e^{(s-1/2)x}\,dx.
\]
For a field of degree $n$ with $r_1$ real embeddings the archimedean
part of the explicit formula is
\begin{align}\label{eq:explicit-arch}
 A={}&n(\gamma+\log(8\pi))+\frac{r_1\pi}{2}
 -n\int_0^\infty\frac{1-F(x)}{2\sinh(x/2)}\,dx
 -r_1\int_0^\infty\frac{1-F(x)}{2\cosh(x/2)}\,dx.
\end{align}
We use the general-signature formula in \cite[Section~4]{miller}.
For $F_0(x)=1/\cosh(x/2)$ the integrals are $\log2$ and $\pi/2-1$.
Replacing $F_0$ by $F$ increases them by at most, respectively,
\[
 \frac1\ell\int_0^\infty\frac{x}{\sinh x}\,dx
 =\frac{\pi^2}{4\ell},\qquad
 \frac1\ell\int_0^\infty\frac{x}{2\cosh^2(x/2)}\,dx
 =\frac{2\log2}{\ell}.
\]
Here the same upper bound includes the tail $x>\ell$ since $1\le x/\ell$.
It follows that
\begin{equation}\label{eq:arch-lower}
 A\ge n\left(\gamma+\log(4\pi)-\frac{\pi^2}{4\ell}\right)
       +r_1\left(1-\frac{2\log2}{\ell}\right).
\end{equation}

If the proposed generating subgroup were proper, its quotient would
have a cyclic quotient of prime order $d$. The ordinary class field
correspondence would give an extension $N/M$ of degree $d$, unramified
at finite primes and split at real places, in which all primes of norm
at most $b$ split completely. Subtract the explicit formula for
$\zeta_M$ from that for $\zeta_N$ and divide by $d-1$.
The pole terms cancel and the discriminant term becomes $\log D_M$.
The remaining zeros belong to the nontrivial abelian Hecke
$L$-functions. Their contributions have nonnegative real part without
assuming the generalized Riemann hypothesis. To see this, for
$|a|<1/2$ the Fourier transform of $\cosh(ax)/\cosh(x/2)$ is
\[
 2\pi\Re\!\left(\operatorname{sech}(\pi(t+ia))\right)\ge0.
\]
The triangle function also has nonnegative Fourier transform, and
multiplication convolves these transforms. At $|a|=1/2$ the ratio is
one. The compact piecewise smooth test is admissible by smoothing;
equivalently its transforms decay quadratically and the zero sum
converges absolutely.

For a prime ideal $\mathfrak p$ of norm $q$ and $t=q^m\le X_*$, the
unsigned prime contribution is
\[
 w(t,m)=\frac{4\log t}{m(t+1)}\left(1-\frac{\log t}{\ell}\right).
\]
Its coefficient is the average of the nontrivial characters of the
cyclic quotient at $\operatorname{Frob}_{\mathfrak p}^m$. It is one
when $t\le b$ and is always at least $-1$. We therefore obtain
\begin{equation}\label{eq:class-bound-contradiction}
 \log D_M\ge A+S_{\mathrm{low}},
\end{equation}
where verified terms are counted positively and every other term may
be counted negatively.

The counter \cert{type13_prym_class/prime_counts.cpp}, in its
\texttt{type18} mode, counts all prime-ideal powers through $X_*$.
It decomposes small and exceptional primes in the maximal order.
For $p>\sqrt{X_*}$ outside the quartic discriminant, only degree-one
prime ideals can contribute. If Frobenius fixes $m$ quartic roots,
it fixes $m(m-1)$ ordered pairs, giving exactly that many degree-one
primes of $M$. This avoids factoring the degree-twelve polynomial at
each such prime.

Counts are binned, retaining the exponent $m$, in intervals
\[
 l=2^e(256+j)/256,\qquad
 u=\min\{X_*,2^e(257+j)/256\},\quad0\le j<256.
\]
A bin with $u\le b$ uses the lower bound
$4\log l(1-\log u/\ell)/(m(u+1))$.
Every other bin, including a bin crossing $b$, uses the negative of
$4\log u(1-\log l/\ell)/(m(l+1))$.
The rational interval computation
\cert{type13_prym_class/check_bound.py}, invoked with \texttt{--type18},
uses $n=r_1=12$ in \eqref{eq:arch-lower} and proves
\[
 A+S_{\mathrm{low}}-\log D_M>0.28.
\]
Logarithms are enclosed using the positive atanh series and a geometric
remainder; each bin contribution is rounded down. Thus the positive
margin contradicts \eqref{eq:class-bound-contradiction} and proves the
bound $b$.

The generation argument uses the ordinary class group throughout.
Passing to its cyclic quotient is what forces real places to split
in the corresponding extension and makes the signatures subtract in
exactly the same proportion as the degrees. The same subtraction
removes the pole term, which would otherwise obstruct a useful bound
with this compact test function. Every prime outside the verified
range is charged its worst allowed sign, so neither the distribution
of those primes among ideal classes nor an unproved splitting density
enters the inequality.

Finally \cert{type18_descent/check_factor_base.gp} applies PARI's exact
prime-ideal relation test through $b$ to the field defined by $B$,
and verifies that every prime in its factor base has norm at most
$1000$. This is an exact reduction of the proven generating set,
not an assumption about a provisional class group or unit index;
see \cite{pari}. Hence $1000$ is a proven generating bound.
With this bound, \cert{type18_descent/rank.m} computes class-group
invariants $[2]$. The custom bound is used only after checking both
the discriminant and a field isomorphism with $M$; other fields retain
Minkowski bounds. This is the proof-bound interface of \cite{magma}.
\end{proof}

\subsection{The rank and the locally soluble quartics}

\begin{proposition}\label{prop:marked-rank}
The curve \eqref{eq:marked-curve} has $\operatorname{rank}E(K)=1$.
\end{proposition}
\begin{proof}
Using Lemma~\ref{lem:type18-class-group}, the exact two-descent in
\cert{type18_descent/rank.m} gives
\[
 \dim_{\F_2}\Sel_2(E/K)=3.
\]
Irreducibility of $q$ shows $E(K)[2]=0$. To distinguish rank one
from rank three, we use two explicit everywhere locally soluble
quartics, rather than infer rank from a point search.

The quartics are $V^2=f_i(u)$, with
$f_i(u)=\sum_{j=0}^4 a_{ij}(h)u^j/219$. For a compact specification,
write a four-tuple $(b_3,b_2,b_1,b_0)$ for
$b_3h^3+b_2h^2+b_1h+b_0$. Their coefficient tuples are
\begingroup\small
\begin{align*}
 a_{14}&=3(326191125074,6862578283530,-228375184609562,359306674735384),\\
 a_{13}&=(-68146848515935,-1146954650973021,42423477924819832,-67064446139806428),\\
 a_{12}&=3(532419837520287,9948662909456332,-349661419343905572,551521283298925706),\\
 a_{11}&=(-17476367657300689,-311397084968974479,11197832892950970757,-17680324780047370650),\\
 a_{10}&=(70099333572212947,1276741669345440462,-45426339226399038625,71690263015899809025),\\[2pt]
 a_{24}&=3(21676496499,-813389539308,8232971883903,-11541580840841),\\
 a_{23}&=(2668272288424,-100136656976148,1013669592704180,-1421062099339992),\\
 a_{22}&=3(13684809418338,-513673043233358,5200629586198914,-7290949365165148),\\
 a_{21}&=(280668222680764,-10541234662342296,106777148298199184,-149707943110710504),\\
 a_{20}&=(719625048806446,-27039042667262238,273964615213522625,-384132484185360732).
\end{align*}
\endgroup
Here the common denominator is part of the specification; no change
of squareclass is being made by clearing it. Each displayed tuple is
interpreted in the fixed power basis of $K$ from
\eqref{eq:quartic-K}, so the quartics can be reconstructed without
choosing embeddings or matching numerical approximations to roots.

The exact certificate \cert{type18_descent/pairing.m} checks that each
quartic is squarefree and everywhere locally soluble. It also checks
the Jacobian identifications: for
$f=au^4+bu^3+cu^2+du+e$, put
\[
 I=12ae-3bd+c^2,\qquad
 J=72ace+9bcd-27ad^2-27b^2e-2c^3.
\]
The curve with coefficients $[0,0,0,-27I,-27J]$ is verified to be
$K$-isomorphic to $E$ for both quartics. The same certificate evaluates
their Cassels--Tate pairing as $1\in\F_2$, in Magma's binary convention.
Thus their images generate a subspace of dimension at least two in
$\Sha(E/K)[2]$: an alternating pairing nonzero on two vectors makes
them independent. The Kummer sequence
\[
 0\longrightarrow E(K)/2E(K)\longrightarrow\Sel_2(E/K)
 \longrightarrow\Sha(E/K)[2]\longrightarrow0
\]
therefore gives rank at most one. This does not require finiteness of
$\Sha$. The exact calculation in Lemma~\ref{lem:dyadic-direction}
shows that $P_0$ has infinite order, supplying the reverse inequality.
\end{proof}

\subsection{Excluding the entire dyadic disk}

\begin{lemma}\label{lem:dyadic-direction}
There is a unique place $v$ of $K$ above $2$, with ramification degree
one and residue degree four. In $L=K_v$, normalize $v(2)=1$.
For $T=24P_0$ on the displayed model,
\[
 v(X(T))=-6,\qquad v(Y(T))=-9,\qquad
 \overline{\frac{-X(T)}{8Y(T)}}=\bar h^3+\bar h\notin\F_2.
\]
The point $P_0$ has infinite order.
\end{lemma}
\begin{proof}
These are exact number-field computations in
\cert{type18_descent/logarithm.gp}. The certificate also verifies that
$k$ is a dyadic unit and that the given point lies on $E$.
Modulo two the defining polynomial of $h$ is
$h^4+h^3+h^2+h+1$, irreducible over $\F_2$; its root has degree four.
In particular $\bar h^3+\bar h$ is neither zero nor one.

The model is integral at $v$. Its invariant differential in the
formal parameter $z=-X/Y$ has an integral power series, so
\[
 \log_E(z)=z+\sum_{i\ge2}b_i z^i/i,\qquad b_i\in\mathcal O_L.
\]
If $v(z)\ge2$, then
$iv(z)-v_2(i)>v(z)$ for every $i\ge2$. Thus the series converges,
preserves valuation, and preserves the leading residue.
For $T$ its parameter has valuation three and nonzero logarithm,
so $T$, and hence $P_0$, has infinite order.
\end{proof}

By Proposition~\ref{prop:marked-rank}, every global elliptic logarithm
lies on the $\Q_2$-line generated by $\log_E(T)$. More explicitly,
extend the logarithm from the subgroup $v(z)\ge2$ to $E(L)$ by
multiplication by an integer taking a point into that subgroup, and
then divide by that integer. The subgroup has finite index, so this
extension is well defined and kills torsion. For any $P\in E(K)$,
rank one gives an integral relation $mP=nT+R$ with $m\ne0$ and $R$
torsion. It follows that $\log_E(P)=(n/m)\log_E(T)$.
Every nonzero element of this line has normalized leading residue
$\bar h^3+\bar h$, since a unit of $\Q_2$ reduces to one.

We now compare this direction with the chord through $P_0$. This
comparison uses a convergent logarithm on the whole admissible disk,
not a finite expansion or a search to a specified rational height.

\begin{proposition}\label{prop:marked-isolation}
A rational $S\equiv1\pmod8$ satisfying
\eqref{eq:five-squareclasses} for $a=25$ and with $H(S)$ a nonzero
rational square must equal $25$.
\end{proposition}
\begin{proof}
Suppose $S\ne25$ and set $j=v(S-25)\ge3$.
Since $q(25)=189^2/k$ is a unit and $q$ has integral dyadic
coefficients, $q(S)/q(25)\in1+8\mathcal O_L$.
The convergent binomial series supplies a square root in
$1+4\mathcal O_L$. Choose the sign of $Y$ in \eqref{eq:marked-curve}
so that $Y/(189k)\equiv1\pmod4$. Then
\[
 Q=P(S)-P_0,\qquad m=\frac{Y+189k}{k(S-25)},\qquad v(m)=1-j.
\]
The addition formulas give
\[
 X(Q)=m^2-kc_2-kS-25k,\qquad
 Y(Q)=-Y+m\bigl(kS-X(Q)\bigr).
\]
All original coordinates and coefficients are integral, whereas
$v(m)\le-2$. Consequently
\[
 X(Q)/m^2\equiv1\pmod2,\qquad
 Y(Q)/(-m^3)\equiv1\pmod2.
\]
In particular $Q$ lies in the formal subgroup with
$v(z(Q))=j-1\ge2$, and the leading residue of $z(Q)$ equals that
of $1/m$. But
\[
 \frac{1/m}{2^{j-1}}
 =\frac{(S-25)/2^j}{189}\,
   \frac{2(189k)}{Y+189k}\equiv1\pmod2.
\]
The first factor is a dyadic rational unit, and the second is one
modulo two by the sign choice. The logarithm therefore has normalized
leading residue one and is nonzero. This contradicts the global
logarithm line, whose nonzero elements have normalized residue
$\bar h^3+\bar h\notin\F_2$. Hence $S=25$.
\end{proof}

We can now finish the progression rigidity statement. The denominator
lemma and Proposition~\ref{prop:five-classes} provide exactly the vertex
required by Proposition~\ref{prop:marked-isolation}. Thus
$U=25H(25)=945^2$, and \eqref{eq:split-quintic} gives
$\{y_i^2\}=\mathcal A$ with multiplicity.

Finally, translate and scale a nonconstant ten-term arithmetic
progression to
\[\{\pm1,\pm3,\pm5,\pm7,\pm9\}.\]
If a rational
multiset has the same first nine power sums, Newton's identities make
its monic degree-ten root polynomial differ from the progression
polynomial by a constant. It is therefore even. Its rational roots
can be arranged in opposite pairs with multiplicity, including zero
pairs if present. The first four even moments of one root from each
pair are precisely \eqref{eq:four-even-moments}. The preceding argument
recovers the five seed squares, so the full ten-term multiset is the
progression. This proves Theorem~\ref{thm:progression} and supplies the
progression case needed for the shift classification.
